\begin{helper}
For every real number \(\rho>0\), there is an integer \(s_0\ge4\) such that
for every integer \(s\ge s_0\),
\[
  (1-\rho)2^{2s-3}
  \le
  2^{2s-3}-2^{s-1}-2^{s-2}+1.
\]
\end{helper}
\begin{proof}
Choose \(s_0\ge4\) so large that \(2^{s-1}\ge 3/\rho\) for every
\(s\ge s_0\).  This is possible because \(2^n\to\infty\).  Fix
\(s\ge s_0\).  Then \(s\ge4\), so all displayed natural-number subtractions
agree with the corresponding integer exponents.

The two lower-order terms satisfy
\[
  2^{s-1}+2^{s-2}
  =3\cdot 2^{s-2}.
\]
Also
\[
  2^{2s-3}=2^{s-1}2^{s-2}.
\]
Since \(2^{s-1}\ge3/\rho\), multiplying by the nonnegative factor
\(\rho\,2^{s-2}\) gives
\[
  2^{s-1}+2^{s-2}
  =3\cdot2^{s-2}
  \le \rho\,2^{s-1}2^{s-2}
  =\rho\,2^{2s-3}.
\]
Therefore
\[
\begin{aligned}
  2^{2s-3}-2^{s-1}-2^{s-2}+1
  &\ge 2^{2s-3}-\bigl(2^{s-1}+2^{s-2}\bigr)\\
  &\ge 2^{2s-3}-\rho\,2^{2s-3}\\
  &=(1-\rho)2^{2s-3}.
\end{aligned}
\]
This proves the claimed eventual lower bound for the F2 vertex count.
\end{proof}
